\scanpage{16}
\begin{flushright}
Massy le 3.11.1970\\
\end{flushright}
Cher Verdier,\\

Ne m'étant guère occupé de Math depuis trois mois, je suis un peu perdu pour SGA 4. Si je me rappelle bien, tu as rédigé ou es en train de rédiger les parties suivantes, qui sont exactement ce qui manque pour que SGA 4 soit complet:
\diagram{\(\begin{array}{ll}
\text{\deleted{IV,}} & \text{\deleted{par.9 et suivants}}\\
\text{V} & \\
\text{VI} & \text{par.5 et suivants}
\end{array}\)}

Si je me rappelle bien, \unsure{\deleted{IV}VI|\deleted{VI}IV} était en fait terminé d'être rédigé, \unsure{\(\mathbb{il}\)|il} fallait seulement y apporter quelques modifications dont on avait discuté avant ton départ. De plus, je pense que tu as avec toi l'exposé
\[
\text{VI B}
\]
de Saint Donat, et je t'envoie également l'Appendice à XVII du même, dont j'avais apparemment lu les premières pages. P\(\mathbb{o}\)urrais tu me le renvoyer (ou le renvoyer à St Donat) avec tes annotations et commentaires ?

Tu dois avoir un\deleted{e} exemplaire du tirage de XVII; le XVIII n'a pas été tapé sur \unsure{\(\mathbb{S}\)|\deleted{s}S|\deleted{S}s}tencils, mais directement sur papier pour offset; \unsure{\deleted{il}la|\deleted{la}il} frappe est terminée, et Deligne est censé la corriger.

Ecris-moi stp où tu en es avec ta part de rédaction, et avec la lecture de VI B . \deleted{S\unsure{\^{r}|\^{t}|î}inger} Pour des raisons techniques, ce serait bien agréable si les textes V,VI,\(\mathrm{VI}_{\mathrm{B}}\) pouvaient être disponibles dans les deux, trois mois qui viennent. Dis-moi en tous cas si tu as l'intention de terminer, et si oui, quand tu penses avoir terminé.

Bien cordialement\\
% archive sheet number: 15
